\section{The number of primary identifiers in a layer}
\label{sec:area_law_primary_count}

Retain the origin, layer and pitch interiors from
Definition~\ref{def:area_law_primary_regions}. Write $I_k$ for the
layer-cell index set, $r=2^k$, $t=2^\ell$, $s=2^p$, and
$o'=o+t(a,b)$. The count below is the primary count in
\cite[Section~11]{OpenAI2026AreaLaw}.

% Source: 10-geometry.tex, lines 212--218 and 237--238.
% The entries count the actual primary interiors in a single layer;
% they do not assert the remaining two-family or repair properties.

\begin{definition}[Primary pitch indices]
    \label{def:area_law_primary_pitch_indices}
    \lean{TNLean.PEPS.AreaLaw.Geometry.primaryPitchIndices}
    \leanok
    \uses{def:area_law_primary_regions,def:area_law_dyadic_cells}
    For every $z\in I_k$, let $L_z$ and $H_z$ be the lower and upper
    corners of the closed cell of side $r$. For each coordinate $i$,
    define
    \begin{align}
        q_{z,i}^- & =\left\lfloor\frac{L_{z,i}-o'_i}{s}\right\rfloor, \\
        q_{z,i}^+ & =\left\lfloor\frac{H_{z,i}-o'_i}{s}\right\rfloor.
    \end{align}
    Let $B_z$ be the integer rectangle
    $[q_{z,1}^-,q_{z,1}^+]\times[q_{z,2}^-,q_{z,2}^+]$.
    The finite set of primary pitch indices is
    \begin{align}
        \mathcal P_k
         & =\left\{j\in\bigcup_{z\in I_k}B_z:
        D_k\cap U_j\ne\varnothing\right\}.
    \end{align}
\end{definition}

\begin{theorem}[Exact primary membership]
    \label{thm:area_law_primary_pitch_membership}
    \lean{TNLean.PEPS.AreaLaw.Geometry.mem_primaryPitchIndices_iff}
    \leanok
    \uses{def:area_law_primary_pitch_indices}
    For every choice of the nonnegative scale indices and every
    $j\in\mathbb Z^2$,
    \begin{align}
        j\in\mathcal P_k & \quad\Longleftrightarrow\quad
        D_k\cap U_j\ne\varnothing.
    \end{align}
\end{theorem}
\begin{proof}\leanok
    \uses{thm:area_law_dyadic_nesting}
    Only the reverse implication needs a proof. A point $x$ in the
    intersection belongs to a cell $C_{k,z}$ with $z\in I_k$.
    Each coordinate lies between those of $L_z$ and $H_z$, so
    monotonicity of the floor gives
    $q_{z,i}^-\le\lfloor(x_i-o'_i)/s\rfloor\le q_{z,i}^+$.
    The open pitch inequalities imply
    $j_i<(x_i-o'_i)/s<j_i+1$; hence the middle integer is $j_i$.
    Thus $j\in B_z$, and the defining filter retains it.
\end{proof}

\begin{theorem}[Primary count from the layer-cell count]
    \label{thm:area_law_primary_pitch_count}
    \lean{TNLean.PEPS.AreaLaw.Geometry.card_primaryPitchIndices_le}
    \leanok
    \uses{def:area_law_primary_pitch_indices}
    If $k\le p$, then
    \begin{align}
        |\mathcal P_k| & \le4|I_k|.
    \end{align}
\end{theorem}
\begin{proof}\leanok
    \uses{def:area_law_dyadic_cells}
    The scale ordering gives $r\le s$. In each coordinate,
    $H_{z,i}=L_{z,i}+r$ and therefore
    \begin{align}
        q_{z,i}^- & \le q_{z,i}^+\le q_{z,i}^-+1.
    \end{align}
    The upper bound follows by comparing $(H_{z,i}-o'_i)/s$ with
    $(L_{z,i}-o'_i)/s+1$ and using
    $\lfloor u+1\rfloor=\lfloor u\rfloor+1$.
    Thus $|B_z|\le4$. A finite union has cardinality at most the sum
    of the cardinalities of its members, and the nonempty-intersection
    filter can only decrease the cardinality. This proves the bound.
\end{proof}

\begin{theorem}[Primary count from the cut boundary]
    \label{thm:area_law_primary_boundary_count}
    \lean{TNLean.PEPS.AreaLaw.Geometry.card_primaryPitchIndices_boundary_le}
    \leanok
    \uses{def:area_law_primary_pitch_indices,def:area_law_hamiltonian}
    Take $Z$ to be the endpoint set of the crossing edges of a cut $A$
    in a finite domain $\Lambda$. If $k\le p$, then
    \begin{align}
        |\mathcal P_k| & \le32(2C_0+1)^2|\partial_\Lambda A|.
    \end{align}
    The bound is uniform in the domain, cut, origin and selected shifts.
\end{theorem}
\begin{proof}\leanok
    \uses{thm:area_law_primary_pitch_count,thm:area_law_dyadic_layer_count}
    Combine $|\mathcal P_k|\le4|I_k|$ with
    $|I_k|\le8(2C_0+1)^2|\partial_\Lambda A|$.
\end{proof}
